\documentclass{situationist-preprint}

\title{Resentment and Forgiveness}
\msid{ACT-09-b}
\msfamily{ACT-09}
\msdiscipline{Action, Agency and Free Will}
\mstype{Consequence}
\mskeywords{forgiveness; resentment; repudiation; protest; remorse; condonation; Hieronymi}

\begin{document}
\maketitle

\begin{abstract}
If resentment is a fitting response to wrongdoing, as ACT-09-a argues, forgiveness seems to require either that the victim cease to represent the wrong correctly or that she give up a fitting attitude for no reason of the right kind. We derive an account of forgiveness from the fittingness of resentment and its separation from the norms of expression. Resentment includes a disposition to protest the wrong and to demand its acknowledgment; when the wrongdoer acknowledges and repudiates the wrong, the demand is met and the reasons for protest lapse. Forgiveness so warranted neither excuses nor condones, leaves the wrongdoer's remorse fitting, and leaves her duties of amends in place.
\end{abstract}

\section{The puzzle}

Forgiveness is commonly understood as the forswearing of resentment \citep{Butler1726,MurphyHampton1988}. It must be distinguished from excusing, which judges that there was no culpable wrong; from justifying, which judges that there was no wrong; from condoning, which accepts the wrong as tolerable; and from forgetting. \citet{Hieronymi2001} put the difficulty sharply: if resentment is a fitting response to a culpable wrong, and forgiveness neither revises the judgment that the wrong was culpable nor forgets it, it is unclear how the victim can give up resentment without either becoming irrational or compromising her judgment. Theorem~3.2 of ACT-09-a makes the puzzle pressing, since it holds that the victim's resentment of a culpable wrong is fitting.

\section{The derivation}

We use Definition~2.1 of ACT-09-a: resentment presents an action as a wrong done to the subject that expresses the agent's ill will or culpable disregard, together with a disposition to protest the wrong and to demand acknowledgment of it.

\begin{definition}[Forgiveness]\label{d:forgive}
To \emph{forgive} a wrongdoer for a wrong is to cease to resent her for it, and to cease to protest it and demand its acknowledgment, without judging that it was not a wrong, that it was not culpable, or that it was acceptable.
\end{definition}

\begin{proposition}\label{p:lapse}
When the wrongdoer acknowledges the wrong and repudiates it, the reasons for continued protest lapse, and forgiveness is warranted without compromise.
\end{proposition}

\begin{proof}
By Lemma~3.1 of ACT-09-a, resentment represents the wrong as expressing the wrongdoer's ill will or culpable disregard. A wrong done in that spirit carries a claim that the victim may be treated so; resentment's protest contests the claim, and its demand for acknowledgment asks the wrongdoer to withdraw it \citep{Hieronymi2001}. When the wrongdoer acknowledges the wrong and repudiates it, she withdraws the claim in her own person: the wrong no longer expresses her present will toward the victim. The demand built into resentment has then been met, and protest has nothing left to contest. Protest is an action, and by Proposition~4.1 of ACT-09-a whether to engage in it depends on practical reasons; its point having been achieved, those reasons lapse. The victim who forgives in these circumstances revises none of the judgments listed in Definition~\ref{d:forgive}: the wrong was a wrong, was culpable, and was unacceptable. So the forgiveness is uncompromising.
\end{proof}

The proposition does not say that the resentment felt at the time was unfitting, or that it becomes unfitting with respect to the wrong as it was done. Theorem~3.2 of ACT-09-a concerns the response to the wrong as an expression of the wrongdoer's will at the time, and that remains a fact. What changes is that continued resentment would present the wrongdoer as still standing behind the wrong, which after her repudiation she does not.

\begin{proposition}\label{p:remains}
Forgiveness does not make the wrongdoer's remorse unfitting, and does not release her from duties of acknowledgment and amends.
\end{proposition}

\begin{proof}
Remorse represents that one has, by one's own agency, infringed a serious claim of another (ETH-27-a). The victim's forgiveness alters her own attitudes and conduct; it does not alter the fact that the claim was infringed, and so does not alter what remorse represents. The duties of acknowledgment and amends are grounded in the same infringement, and survive it in the same way.
\end{proof}

This agrees with ETH-27, where remorse is fitting even when the act was right, since it responds to the infringement of a claim rather than to the wrongness of the act. Forgiveness and remorse are therefore not rivals: the victim's forgiveness answers the wrongdoer's repudiation, and the wrongdoer's remorse continues to answer what she did.

\section{Objections and replies}

\begin{objection}[Gift forgiveness]
Many forgive the unrepentant, and seem admirable for it. On the account given, the reasons for protest have not lapsed, so such forgiveness looks like condonation.
\end{objection}

\begin{reply}
Proposition~\ref{p:lapse} states when forgiveness is warranted by the lapse of the reasons for protest; it does not say that forgiveness is permissible only then. Ceasing to protest is an action, and by Proposition~4.1 of ACT-09-a it may be supported by reasons that do not bear on the fittingness of resentment: the victim's peace, the value of a relationship, the futility of further protest. Forgiveness of the unrepentant is permitted for such reasons, but not owed. It does not condone, provided the victim continues to judge the wrong a wrong and unacceptable; what it gives up is the demand, not the judgment.
\end{reply}

\begin{objection}[Transaction]
\citet{Nussbaum2016} objects that forgiveness conditioned on apology turns a moral relation into an exchange, in which the wrongdoer pays in humiliation for the release of the victim's anger.
\end{objection}

\begin{reply}
The repudiation that warrants forgiveness is not a payment and need not humiliate. It is the wrongdoer's own withdrawal of the claim her wrong made, which is what the victim's protest asked for in the first place. The condition is not imposed from outside the relation; it is what resentment, as an attitude of protest, was always directed at \citep{Strawson1962}.
\end{reply}

\section{Conclusion}

The fittingness of resentment and the separation of the emotion from the norms of its expression yield an account of forgiveness. Resentment protests a wrong and demands acknowledgment; when the wrongdoer acknowledges and repudiates the wrong, the demand is met and the reasons for protest lapse, so that the victim may forgive without excusing, justifying or condoning. The wrongdoer's remorse remains fitting, and her duties of amends remain in force.

\begin{thebibliography}{9}\small
\bibitem[Butler(1726)]{Butler1726} Butler, J. (1726). \emph{Fifteen Sermons Preached at the Rolls Chapel}. Sermons VIII and IX. London.
\bibitem[Hieronymi(2001)]{Hieronymi2001} Hieronymi, P. (2001). Articulating an uncompromising forgiveness. \emph{Philosophy and Phenomenological Research}, 62(3), 529--555.
\bibitem[Murphy and Hampton(1988)]{MurphyHampton1988} Murphy, J.~G. and Hampton, J. (1988). \emph{Forgiveness and Mercy}. Cambridge University Press.
\bibitem[Nussbaum(2016)]{Nussbaum2016} Nussbaum, M.~C. (2016). \emph{Anger and Forgiveness: Resentment, Generosity, Justice}. Oxford University Press.
\bibitem[Strawson(1962)]{Strawson1962} Strawson, P.~F. (1962). Freedom and resentment. \emph{Proceedings of the British Academy}, 48, 1--25.
\end{thebibliography}

\end{document}
